\chapter{Środowisko i architektura}
\label{chap:architektura}

Architektura środowiska laboratoryjnego \textbf{VulnBox} została zaprojektowana zgodnie z zasadą separacji warstw sieciowych oraz stopniowej eskalacji kontekstów wykonawczych w systemie Linux \cite{nist800115}. Rozdział ten przedstawia ogólny układ topologii sieciowej w środowisku wirtualizacji oraz koncepcyjny przepływ danych pomiędzy usługami tworzącymi ścieżkę weryfikacyjną od stanu zerowej wiedzy (\texttt{0-knowledge}) do uprawnień administratora systemu (\texttt{root}). Szczegółowa zawartość plików konfiguracyjnych oraz analiza kodu źródłowego poszczególnych usług zostały omówione w rozdziale \ref{chap:konfiguracja}.

\section{Topologia sieciowa środowiska laboratoryjnego}
\label{sec:topologia}

Środowisko uruchomieniowe bazuje na wirtualizatorze Oracle VirtualBox, w którym maszyna docelowa \texttt{VulnBox} oraz stacja analityczna komunikują się za pośrednictwem wydzielonego, izolowanego segmentu sieci wirtualnej (\emph{Host-Only Ethernet Adapter}, \texttt{vboxnet0}, podsieć \texttt{192.168.56.0/24}). Takie rozwiązanie gwarantuje pełną izolację ruchu laboratoryjnego od sieci zewnętrznej przy jednoczesnym zachowaniu realistycznych warunków komunikacji TCP/IP.

\TikzFigure[0.95]{topologia_sieci.tex}{Ogólna topologia sieciowa środowiska laboratoryjnego VulnBox}{fig:topologia}

Jak przedstawiono na rysunku \ref{fig:topologia}, komunikacja z hosta analitycznego przechodzi przez wirtualny przełącznik sieci \texttt{vboxnet0} bezpośrednio do interfejsu sieciowego maszyny \texttt{VulnBox}, chronionego przez zaporę pakietową UFW oraz demona Fail2Ban. Na etapie budowy i instalacji pakietów systemowych maszyna może tymczasowo korzystać z pomocniczego interfejsu NAT, który przed finalnym eksportem do pliku \texttt{VulnBox.ova} zostaje wyłączony, pozostawiając wyłącznie kartę sieci laboratoryjnej skonfigurowaną do automatycznego pobierania adresu IP z serwera DHCP.

\section{Architektura usług i teoretyczna ścieżka weryfikacji}
\label{sec:architektura-uslug}

Wewnętrzna architektura maszyny \texttt{VulnBox} obejmuje cztery powiązane ze sobą warstwy usług działające w odrębnych kontekstach kont systemowych (\texttt{nobody}, \texttt{www-data}, \texttt{dev} oraz \texttt{root}). Zamiast eksponować pojedynczą, trywialną usługę sieciową, projekt modeluje realistyczny scenariusz, w którym drobne nieprawidłowości konfiguracyjne w kolejnych komponentach łączą się w spójny łańcuch zależności.

\TikzFigure[0.98]{architektura_warstw.tex}{Powiązanie usług oraz teoretyczna ścieżka przejścia od stanu 0-knowledge do konta root}{fig:architektura}

Na rysunku \ref{fig:architektura} zobrazowano pięć kolejnych kroków składających się na teoretyczną ścieżkę weryfikacji środowiska:
\begin{enumerate}
    \item \textbf{Stan początkowy (\texttt{0-knowledge})} - osoba weryfikująca maszynę dysponuje wyłącznie dostępem do segmentu sieciowego \texttt{vboxnet0}, nie posiadając żadnych poświadczeń ani informacji o strukturze katalogów na serwerze \texttt{vulnbox}.
    \item \textbf{Etap 1 - Usługa plikowa SMB (\texttt{:445/tcp}, kontekst \texttt{nobody})} - serwer Samba udostępnia zasób gościnny \texttt{[public]} niewymagający uwierzytelniania. Pozostawione w nim notatki wdrożeniowe ujawniają niestandardowe ścieżki instalacyjne portalu diagnostycznego WWW oraz lokalnego serwisu zarządzającego.
    \item \textbf{Etap 2 - Portal diagnostyczny WWW (\texttt{:8080/tcp}, kontekst \texttt{www-data})} - aplikacja PHP obsługiwana przez serwer Nginx zawiera celowo zaprojektowaną, jednoprzebiegową normalizację ścieżek plików logów (\texttt{?log=}) \cite{cwe22} oraz moduł mostka diagnostycznego (\texttt{?proxy\_path=}) komunikującego się z interfejsem pętli zwrotnej (\texttt{127.0.0.1}). Pozwala to na odczyt statycznego tokenu autoryzacyjnego z pliku konfiguracyjnego wewnętrznego API.
    \item \textbf{Etap 3 - Wewnętrzne API zarządzające (\texttt{127.0.0.1:5000}, kontekst \texttt{dev})} - mikroserwis Flask nasłuchuje wyłącznie na interfejsie lokalnym \texttt{loopback} (niedostępnym bezpośrednio z sieci zewnętrznej) i działa z uprawnieniami konta \texttt{dev}. Wykorzystując pozyskany w etapie 2 token \texttt{X-Internal-Token} oraz mostek HTTP portalu WWW, możliwe jest wywołanie punktu końcowego rejestrującego klucz publiczny SSH w pliku \texttt{/home/dev/.ssh/authorized\_keys}.
    \item \textbf{Etap 4 i Etap 5 - Dostęp lokalny (\texttt{dev}) oraz przejście do konta \texttt{root}} - po nawiązaniu sesji SSH jako użytkownik \texttt{dev} i odczycie flagi \texttt{/home/dev/user.txt}, architektura systemu przewiduje dwie niezależne ścieżki uzyskania uprawnień \texttt{uid=0(root)} i odczytu flagi \texttt{/root/root.txt}:
    \begin{itemize}
        \item \emph{Ścieżka A (zadanie cykliczne \texttt{cron})} - wykorzystanie uprawnień zapisu grupy \texttt{dev} do skryptu pomocniczego \texttt{/opt/admin/hooks/pre\_backup.sh} wywoływanego co minutę przez zadanie \texttt{/etc/cron.d/vulnbox-maintenance} w kontekście \texttt{root},
        \item \emph{Ścieżka B (delegacja \texttt{sudoers.d})} - wykorzystanie reguły \texttt{NOPASSWD: SETENV} dla polecenia \texttt{/usr/bin/python3 /opt/admin/maintenance.py}, które importuje moduł \texttt{status\_helper.py} z katalogu \texttt{/opt/admin} oraz respektuje zmienną środowiskową \texttt{PYTHONPATH}.
    \end{itemize}
\end{enumerate}
